\section{Conclusion}
\label{sec:conclusion}

We tested whether tokenizer fertility explains zero-shot cross-lingual transfer loss in \mbert and \xlmr across three tasks and up to 71 target languages.
Fertility explains 0--27\% of the between-language variance in transfer drop, not 40--60\%.
Syntactic distance outperforms it in every setting we tested, the 2.5$\times$ threshold is not a reliable failure boundary, and fertility's predictive value beyond pretraining size and syntactic distance is small.
Fragmentation does have a causal cost when it is unfamiliar to the model, and a tokens-per-parallel-sentence measure shows a stronger association that survives controls on one task.
The key takeaway is that how finely a tokenizer splits a language tells us little, by itself, about how well a model will transfer to it.

Several follow-ups would sharpen these conclusions.
The most valuable is to train the same architecture on the same corpus with different vocabularies; this is the experiment that isolates the tokenizer, and nothing here does.
Testing the parallel token ratio on other parallel benchmarks would show whether the \sib result is general or an artifact of measuring fertility on the task's own text distribution.
Replacing raw fertility with a measure of \emph{excess} fragmentation, for example relative to a monolingual tokenizer~\cite{rust2021good} or to morpheme counts, would separate the two components we propose in \secref{sec:discussion}.
Finally, the analysis should be repeated with full-data \xnli and with larger and decoder-only models.
